\IfClozeTF{\clearpage}{\medskip}
\subsection{Quadratische Gleichungen}
\label{subsec:fundamentals_quadratic_equations}

Eine quadratische Gleichung in einer Variablen hat die Form
\[
  ax^2+bx+c=0, \qquad a\neq 0.
\]
\noindent\begin{minipage}{\linewidth}
  Mit Hilfe der quadratischen Ergänzung können wir die Gleichung nach $x$ auflösen:
  \begin{ClozeArea}[34]{24}
    \[
      \begin{array}{r@{\quad}l@{\qquad}l}
        & ax^2+bx+c=0
          & \TransformOp{-c} \\[1.1em]
        \Longleftrightarrow & ax^2+bx=-c
          & \TransformOp{:a \quad (a\neq0)} \\[1.1em]
        \Longleftrightarrow & x^2+\dfrac{b}{a}x=-\dfrac{c}{a}
          & \TransformOp{+\dfrac{b^2}{4a^2}
              \qquad\text{\enquote{quadratische Ergänzung}}} \\[1.1em]
        \Longleftrightarrow
          & x^2+\dfrac{b}{a}x{\color{THMBlue}+\dfrac{b^2}{4a^2}}
            =-\dfrac{c}{a}{\color{THMBlue}+\dfrac{b^2}{4a^2}}
          & \\[1.1em]
        \Longleftrightarrow
          & \left(x+\dfrac{b}{2a}\right)^2
            =\dfrac{b^2-4ac}{4a^2}
          & \TransformOp{\pm\sqrt{\vphantom{b^2}}} \\[1.1em]
        \Longleftrightarrow
          & x+\dfrac{b}{2a}
            =\pm\dfrac{\sqrt{b^2-4ac}}{2a}
          & \TransformOp{-\dfrac{b}{2a}} \\[1.1em]
        \Longleftrightarrow
          & x=-\dfrac{b}{2a}\pm\dfrac{\sqrt{b^2-4ac}}{2a}.
          &
      \end{array}
    \]
  \end{ClozeArea}
\end{minipage}

\begin{TheoremBox}{abc-Formel}{fundamentals_quadratic_formula}
  Die Lösungen der quadratischen Gleichung $ax^2+bx+c=0$ mit $a\neq0$ lauten
    \[
      x_{1,2}=\frac{-b\pm\sqrt{b^2-4ac}}{2a}.
    \]
\end{TheoremBox}

\begin{RemarkBox}[remark title=Diskriminante]
  Die \emph{Diskriminante} $D=b^2-4ac$ bestimmt die Anzahl der reellen
  Lösungen einer quadratischen Gleichung:
  \[
    \begin{array}{ccl}
      D<0 &\Longleftrightarrow& \Cloze{8cm}{\text{keine reelle Lösung}},\\[0.7em]
      D=0 &\Longleftrightarrow& \Cloze{8cm}{\text{eine reelle Lösung}},\\[0.7em]
      D>0 &\Longleftrightarrow& \Cloze{8cm}{\text{zwei reelle Lösungen}}.
    \end{array}
  \]
\end{RemarkBox}

\begin{ExampleBox}{Quadratische Gleichung}{fundamentals_quadratic_equation_example}
  Wir bestimmen, wie viele reelle Lösungen die folgende quadratische Gleichung
  besitzt, und berechnen diese:
  \[
    2x^2-6x+4=0.
  \]
  \begin{ClozeArea}[34]{16}
    Da die Diskriminante
    \[
      D=b^2-4ac=(-6)^2-4\cdot2\cdot4=4>0.
    \]
    positiv ist, besitzt die Gleichung zwei reelle Lösungen. Mit der abc-Formel folgt
    \begin{align*}
      &x_{1,2} = \frac{-b\pm\sqrt{b^2-4ac}}{2a} = \frac{-(-6)\pm\sqrt{(-6)^2-4\cdot2\cdot4}}{2\cdot2} = \frac{6\pm2}{4} \\
      \Longrightarrow \qquad &x_1=1, \qquad x_2=2.
    \end{align*}
  \end{ClozeArea}
  Die Lösungsmenge enthält $\Cloze{1cm}{2}$ Elemente und lautet $\mathcal{L}=\Cloze{5cm}{\{1,2\}}$.
\end{ExampleBox}
